% ថ្នាក់ទី ១១ (វិទ្យាសាស្ត្រសង្គម) · ជំពូកទី ៧ · មេរៀនទី ១ · សៀវភៅ ទំព័រ ២៦៧–២៧៥
\chapter{រង្វាស់ទីតាំងនៃទិន្នន័យ}

\begin{note}[ជំពូកទី ៧ ៖ ស្ថិតិ]
ជំពូកនេះមានមេរៀន ៖ \textbf{១. រង្វាស់ទីតាំងនៃទិន្នន័យ} \quad \textbf{២. រង្វាស់នៃគម្លាត} ។

នៅជំពូកនេះ យើងនឹងសិក្សាពីការបែងចែកសំណុំទិន្នន័យមួយជាភាគរយ និងបម្រែបម្រួលសំណុំទិន្នន័យមួយធៀបនឹងសំណុំទិន្នន័យមួយទៀត ។ រង្វាស់ដែលចាំបាច់សម្រាប់ពិពណ៌នាទិន្នន័យគឺជារង្វាស់ទីតាំង ។ រង្វាស់នេះប្រាប់ពីភាគរយនៃទិន្នន័យដែលមានតម្លៃតិច ឬច្រើនជាងតម្លៃណាមួយជាក់លាក់ ។ ម្យ៉ាងទៀតទិន្នន័យអាចមានភាពប្រែប្រួលតិច ឬច្រើនពីតម្លៃមធ្យម ដែលបណ្ដាលឱ្យបំណែងចែកទិន្នន័យមានពង្រាយខ្លាំង ឬខ្សោយ ។ គេបានបង្កើតរង្វាស់គម្លាតដើម្បីវាស់វែងពីបម្រែបម្រួលទិន្នន័យ ។ ចំពោះការត្រួតពិនិត្យទំនិញ គេពិនិត្យមើលលើតម្លៃគម្លាតនេះ ដែលបានកំណត់ក្នុងស្តង់ដារនៃគុណភាពទំនិញរបស់គេ ។
\end{note}

\begin{objBox}[វត្ថុបំណង]
\begin{enumerate}[1]
  \item គណនាកាទីល ឬភាគបួន ភាគដប់ និងភាគរយនៃទិន្នន័យមិនផ្ដុំជាតារាងប្រេកង់
  \item គណនាកាទីល ឬភាគបួន ភាគដប់ និងភាគរយនៃទិន្នន័យផ្ដុំជាតារាងប្រេកង់ ។
\end{enumerate}
\end{objBox}

\section{រង្វាស់ទីតាំងនៃទិន្នន័យមិនផ្ដុំជាតារាងប្រេកង់}

\begin{example}[ឧទាហរណ៍]
សាលារៀនបើកបរថយន្តមួយកន្លែងទទួលបានសិស្សចំនួន $16$ នាក់ ដែលមានអាយុ $25, 23, 26, 30, 37, 20, 21, 23, 31, 34, 26, 34, 31, 26, 35, 32$ ។ នៅពេលគេរៀបអាយុសិស្សតាមលំដាប់ពីតិចទៅច្រើន រួចហើយគេចង់ដឹងថា ៖ តើ $25\%$ នៃសិស្សទាំងអស់នៅក្រោមអាយុប៉ុន្មានឆ្នាំ? $50\%$ នៃសិស្សទាំងអស់នៅក្រោមអាយុប៉ុន្មានឆ្នាំ? $75\%$ នៃសិស្សទាំងអស់នៅក្រោមអាយុប៉ុន្មានឆ្នាំ?
\end{example}
ចំពោះទិន្នន័យដែលរៀបតាមលំដាប់កើន តម្លៃ $25\%$ នៃទិន្នន័យទាំងអស់ ហៅថា \textbf{កាទីលទី $1$} តាងដោយ $Q_1$ និង $75\%$ នៃទិន្នន័យទាំងអស់ ហៅថា\textbf{កាទីលទី $3$} តាងដោយ $Q_3$ ។ ចំណែក\textbf{កាទីលទី $2$} តាងដោយ $Q_2$ គឺជាមេដ្យាន ជាតម្លៃត្រូវនឹង $50\%$ នៃទិន្នន័យទាំងអស់ ។
\begin{center}
\begin{tikzpicture}[x=0.85cm,y=1.2cm]
\draw (0,0) -- (16,0);
\draw (0,-0.12) -- (0,0.12) (16,-0.12) -- (16,0.12);
\foreach \p/\q in {4/1,8/2,12/3} {\draw (\p,-0.15) -- (\p,0.15); \node[below] at (\p,-0.1) {$Q_{\q}$};}
\foreach \c in {2,6,10,14} \node[above] at (\c,0.05) {\small $25\%$};
\node[left] at (0,0) {\small តម្លៃតូចបំផុត};
\node[right] at (16,0) {\small តម្លៃធំបំផុត};
\foreach \v [count=\i] in {20,21,23,23,25,26,26,26,30,31,31,32,34,34,35,37} \node at ({\i - 0.5},-1) {\small $\v$};
\foreach \p/\q in {4/1,8/2,12/3} {\draw[densely dashed] (\p,-0.75) -- (\p,-1.25); \node at (\p,-1.5) {\small $Q_{\q}$};}
\end{tikzpicture}
\end{center}
គេសង្កេតឃើញថា ចំនួនទិន្នន័យទាំងអស់ស្មើនឹង $16$ ដូច្នេះ ៖
\begin{itemize}
  \item $Q_1 =$ តម្លៃតួទី $\dfrac{1}{4}(16 + 1) =$ តម្លៃតួទី $4.25$ ។ ដោយតួទី $4.25$ នៅចន្លោះតួទី $4$ និងតួទី $5$ គេបាន $Q_1 = \dfrac{23 + 25}{2} = 24$ ។
  \item $Q_2 =$ តម្លៃតួទី $\dfrac{1}{2}(16 + 1) =$ តម្លៃតួទី $8.5$ ។ ដោយតួទី $8.5$ នៅចន្លោះតួទី $8$ និងតួទី $9$ គេបាន $Q_2 = \dfrac{26 + 30}{2} = 28$ ។
  \item $Q_3 =$ តម្លៃតួទី $\dfrac{3}{4}(16 + 1) =$ តម្លៃតួទី $12.75$ ។ ដោយតួទី $12.75$ នៅចន្លោះតួទី $12$ និងតួទី $13$ គេបាន $Q_3 = \dfrac{32 + 34}{2} = 33$ ។
\end{itemize}
\textbf{បំណកស្រាយ ៖} យើងឃើញថា $25\%$ នៃសិស្សទាំងអស់មានអាយុតិចជាង $24$ ឆ្នាំ $50\%$ នៃសិស្សទាំងអស់មានអាយុតិចជាង $28$ ឆ្នាំ ហើយ $75\%$ នៃសិស្សទាំងអស់មានអាយុតិចជាង $33$ ឆ្នាំ ។

\begin{formulaBox}[ជាទូទៅ]
បើទិន្នន័យដែលបានរៀបតាមលំដាប់មានចំនួនតួស្មើនឹង $n$ នោះគេបាន ៖
\[
\begin{aligned}
&\text{កាទីលទី } 1 : && Q_1 = \text{តម្លៃតួទី } \frac{1}{4}(n + 1)\\
&\text{កាទីលទី } 2 : && Q_2 = \text{តម្លៃតួទី } \frac{2}{4}(n + 1) = \text{តម្លៃតួទី } \frac{1}{2}(n + 1)\\
&\text{កាទីលទី } 3 : && Q_3 = \text{តម្លៃតួទី } \frac{3}{4}(n + 1)
\end{aligned}
\]
\end{formulaBox}

\begin{note}[វិធាន]
\begin{itemize}
  \item មុននឹងគណនាកាទីល ត្រូវរៀបទិន្នន័យតាមលំដាប់ពីតូចទៅធំ ។
  \item បើទីតាំងនៃ $Q$ ស្ថិតនៅចំទិន្នន័យណា នោះ $Q$ ស្មើនឹងតម្លៃទិន្នន័យនោះ ។
  \item បើទីតាំងនៃ $Q$ ស្ថិតនៅចំចន្លោះទិន្នន័យពីរ នោះតម្លៃ $Q$ ជាមធ្យមនៃទិន្នន័យទាំងពីរ ។
\end{itemize}
\end{note}

ដូចគ្នានេះដែរ \textbf{ភាគដប់}ចែកទិន្នន័យជាដប់ចំណែកស្មើៗគ្នា ។
\begin{center}
\begin{tikzpicture}[x=1.1cm,y=1.2cm]
\draw (0,0) -- (10,0);
\foreach \i in {0,...,10} \draw (\i,-0.1) -- (\i,0.1);
\foreach \i in {0,...,9} \node[above] at ({\i + 0.5},0.05) {\scriptsize $10\%$};
\foreach \i in {1,...,5} \node[below] at (\i,-0.05) {\small $D_{\i}$};
\node[below] at (9,-0.05) {\small $D_9$};
\node[below] at (7,-0.05) {\small $\cdots$};
\node[below left] at (0,-0.05) {\small អប្បបរមា};
\node[below right] at (10,-0.05) {\small អតិបរមា};
\node[above] at (5,0.55) {\small មេដ្យាន};
\end{tikzpicture}
\end{center}
បើគេតាង $D_i$ ជាទីតាំងនៃភាគដប់ នោះគេបាន $D_1, D_2, D_3, \dots, D_9$ ។
\begin{formulaBox}[រូបមន្តភាគដប់]
\[
\begin{gathered}
D_i = \text{តម្លៃតួទី } \frac{i}{10}(n + 1),\\
i = 1, 2, \dots, 9
\end{gathered}
\]
\end{formulaBox}
\textbf{ភាគរយ}ចែកទិន្នន័យជាមួយរយចំណែកស្មើៗគ្នា ។ បើគេតាង $L_P$ ជាទីតាំងភាគរយ នោះគេបាន $L_1, L_2, \dots, L_{99}$ ។
\begin{formulaBox}[រូបមន្តភាគរយ]
\[
L_P = \text{តម្លៃតួទី } \frac{P}{100}(n + 1)
\]
ដែល $P$ ជាភាគរយដែលគេត្រូវរក ។
\end{formulaBox}

\begin{example}[លំហាត់គំរូ]
ប្រាក់កម្រៃជើងសាប្រចាំខែរបស់បុគ្គលិកក្នុងក្រុមហ៊ុនលក់កុំព្យូទ័រមួយ (គិតជាដុល្លារ) មាន $120, 150, 150, 180, 200, 120, 170, 220, 190$ ។ គណនា $Q_1$ និង $D_2$ ។
\ans
រៀបទិន្នន័យតាមលំដាប់ ៖ $120, 120, 150, 150, 170, 180, 190, 200, 220$ ($n = 9$) ។
\begin{itemize}
  \item $Q_1 =$ តម្លៃតួទី $\dfrac{1}{4}(9 + 1) =$ តម្លៃតួទី $2.5$ ។ ដូច្នេះ $Q_1 = \dfrac{120 + 150}{2} = \$135$ មានន័យថា $25\%$ នៃបុគ្គលិកទាំងអស់មានប្រាក់កម្រៃជើងសាតិចជាង $\$135$ ឬ $75\%$ នៃបុគ្គលិកទាំងអស់មានប្រាក់កម្រៃជើងសាច្រើនជាង $\$135$ ។
  \item $D_2 =$ តម្លៃតួទី $\dfrac{2}{10}(9 + 1) =$ តម្លៃតួទី $2$ ។ ដូច្នេះ $D_2 = \$120$ មានន័យថា $80\%$ នៃបុគ្គលិកទាំងអស់មានប្រាក់កម្រៃជើងសាច្រើនជាង $\$120$ ឬ $20\%$ នៃបុគ្គលិកទាំងអស់មានប្រាក់កម្រៃជើងសាតិចជាង $\$120$ ។
\end{itemize}
\end{example}

\begin{pro}
គណនានិងបញ្ជាក់អត្ថន័យ $Q_1, Q_2, Q_3$ និង $L_{25}$ នៃសំណុំទិន្នន័យខាងក្រោម ៖ $4, 6, 18, 25, 9, 16, 22, 5, 20, 4, 8$ ។
\ans
រៀបតាមលំដាប់ ៖ $4, 4, 5, 6, 8, 9, 16, 18, 20, 22, 25$ ($n = 11$) ។ $Q_1 =$ តួទី $3 = 5$, $Q_2 =$ តួទី $6 = 9$, $Q_3 =$ តួទី $9 = 20$, $L_{25} =$ តួទី $\dfrac{25}{100}\times 12 = 3$ គឺ $L_{25} = 5$ ។ មានន័យថា $25\%$ នៃទិន្នន័យតិចជាង $5$, $50\%$ តិចជាង $9$ និង $75\%$ តិចជាង $20$ ។
\end{pro}

\section{រង្វាស់ទីតាំងនៃទិន្នន័យផ្ដុំជាតារាងប្រេកង់}
ដើម្បីកំណត់រង្វាស់ទីតាំងនៃទិន្នន័យផ្ដុំជាតារាងប្រេកង់ គេត្រូវបង្កើតតារាងប្រេកង់កើន ។

\begin{example}[លំហាត់គំរូ]
អ្នកលក់សៀវភៅម្នាក់បានឱ្យដឹងថា សៀវភៅថ្មីរបស់គាត់ទាំងអស់មានចំនួន $30$ ក្បាល ។ ក្នុងនោះមានចំនួន $5$ ក្បាល មួយក្បាលលក់តម្លៃ $\$12$ មានចំនួន $8$ ក្បាល មួយក្បាលលក់តម្លៃ $\$9$ មានចំនួន $6$ ក្បាល មួយក្បាលលក់តម្លៃ $\$7$ មានចំនួន $9$ ក្បាល មួយក្បាលលក់តម្លៃ $\$15$ និងមានចំនួន $2$ ក្បាល មួយក្បាលលក់តម្លៃ $\$20$ ។ គណនា $Q_1$ មេដ្យាននិង $L_{81}$ ។
\ans
បង្កើតតារាងប្រេកង់កើន ៖
\[
\begin{array}{|c|c|c|}\hline
\text{តម្លៃសៀវភៅ } x & \text{ចំនួនសៀវភៅ } f & f^{\uparrow}\\ \hline
\$7 & 6 & 6\\
\$9 & 8 & 14\\
\$12 & 5 & 19\\
\$15 & 9 & 28\\
\$20 & 2 & 30\\ \hline
\end{array}
\]
\begin{itemize}
  \item $Q_1 =$ តម្លៃតួទី $\dfrac{1}{4}(n + 1)$, $n = 30$ ៖ $Q_1 =$ តម្លៃតួទី $\dfrac{1}{4}(30 + 1) =$ តម្លៃតួទី $7.75$ ។ ដូច្នេះ $Q_1 = \$9$ ។
  \item $Q_2 =$ មេដ្យាន $=$ តម្លៃតួទី $\dfrac{1}{2}(n + 1) =$ តម្លៃតួទី $\dfrac{1}{2}(30 + 1) =$ តម្លៃតួទី $15.5$ ។ ដូច្នេះ មេដ្យានស្មើនឹង $\$12$ ។
  \item $L_{81} =$ តម្លៃតួទី $\dfrac{P}{100}(n + 1) =$ តម្លៃតួទី $\dfrac{81}{100}(30 + 1) =$ តម្លៃតួទី $25.11$ ។ ដូច្នេះ $L_{81} = \$15$ ។
\end{itemize}
\end{example}

\begin{pro}
តារាងបំណែងចែកប្រេកង់ខាងក្រោម បង្ហាញពីចំនួនក្មេងក្នុងគ្រួសារនៃ $35$ គ្រួសារដែលគេបានទៅសម្ភាសន៍ក្នុងភូមិមួយ ។ គណនានិងបកស្រាយមេដ្យាន កាទីលទី $1$ $Q_1$ និង $L_{73}$ ។
\[
\begin{array}{|l|cccccc|}\hline
\text{ចំនួនក្មេង } x & 0 & 1 & 2 & 3 & 4 & 5\\ \hline
\text{ចំនួនគ្រួសារ } f & 3 & 5 & 12 & 9 & 4 & 2\\ \hline
\end{array}
\]
\ans
$f^{\uparrow}$ ៖ $3, 8, 20, 29, 33, 35$ ។ មេដ្យាន $=$ តួទី $18 = 2$ ក្មេង, $Q_1 =$ តួទី $9 = 2$ ក្មេង, $L_{73} =$ តួទី $\dfrac{73}{100}\times 36 = 26.28$ គឺ $L_{73} = 3$ ក្មេង ។ មានន័យថា ពាក់កណ្ដាលនៃគ្រួសារមានក្មេងមិនលើសពី $2$ នាក់ និងប្រហែល $73\%$ នៃគ្រួសារមានក្មេងមិនលើសពី $3$ នាក់ ។
\end{pro}

\section{រង្វាស់ទីតាំងនៃទិន្នន័យផ្ដុំជាថ្នាក់}
ដើម្បីកំណត់រង្វាស់ទីតាំងនៃទិន្នន័យផ្ដុំជាថ្នាក់ គេត្រូវបង្កើតតារាងប្រេកង់កើន ។

\begin{example}[លំហាត់គំរូ ១]
ខាងក្រោមជាតារាងបំណែងចែកប្រេកង់នៃពិន្ទុសិស្សចំនួន $40$ នាក់ ដែលគេបានធ្វើតេស្តគណិតវិទ្យា ។
\[
\begin{array}{|c|c|}\hline
\text{ថ្នាក់នៃពិន្ទុ} & \text{ចំនួនសិស្ស}\\ \hline
30 - 40 & 5\\
40 - 50 & 7\\
50 - 60 & 12\\
60 - 70 & 11\\
70 - 80 & 3\\
80 - 90 & 2\\
\hline
\end{array}
\]
គណនា $Q_1$ និង $Q_3$ ។
\ans
ដំបូងត្រូវបង្កើតតារាងប្រេកង់កើន ៖
\[
\begin{array}{|c|c|c|}\hline
\text{ថ្នាក់នៃពិន្ទុ} & \text{ចំនួនសិស្ស} & f^{\uparrow}\\ \hline
30 - 40 & 5 & 5\\
40 - 50 & 7 & 12\\
50 - 60 & 12 & 24\\
60 - 70 & 11 & 35\\
70 - 80 & 3 & 38\\
80 - 90 & 2 & 40\\ \hline
\end{array}
\]
\begin{itemize}
  \item មួយភាគបួននៃប្រេកង់សរុប ឬ $25\%$ នៃទិន្នន័យគឺ $\dfrac{40}{4} = 10$ ។ ក្នុងជួរប្រេកង់កើន $10$ ស្ថិតនៅក្នុងថ្នាក់ $40 - 50$ ។ ដូច្នេះ $Q_1$ ស្ថិតនៅក្នុងថ្នាក់ $40 - 50$ ។ គេអាចកំណត់ទីតាំង $Q_1$ ដោយប្រើមេគុណប្រាប់ទិសនៃបន្ទាត់ $AC$ និង $AB$ ។ ដោយមេគុណប្រាប់ទិសនៃបន្ទាត់ $AC$ និង $AB$ ស្មើគ្នា គេបាន
  \[
  \frac{10 - 5}{Q_1 - 40} = \frac{7}{10} \Rightarrow Q_1 = 40 + \frac{10 - 5}{7}\times 10 \approx 47.14
  \]
  \item បីភាគបួននៃប្រេកង់សរុប ឬ $75\%$ នៃទិន្នន័យស្មើនឹង $\dfrac{3}{4}\times 40 = 30$ ។ ក្នុងជួរប្រេកង់កើន $30$ ស្ថិតនៅក្នុងថ្នាក់ $60 - 70$ ដូច្នេះ $Q_3$ ស្ថិតនៅក្នុងថ្នាក់ $60 - 70$ ។ គេដឹងថា $m_{A'C'} = m_{A'B'}$ គេបាន
  \[
  \frac{30 - 24}{Q_3 - 60} = \frac{11}{10} \Rightarrow Q_3 = 60 + \frac{30 - 24}{11}\times 10 \approx 65.45
  \]
\end{itemize}
\begin{center}
\begin{tikzpicture}[x=0.22cm,y=0.42cm]
\draw[->] (30,0) -- (58,0) node[right] {\small ពិន្ទុ};
\draw[->] (30,0) -- (30,14) node[left] {$f^{\uparrow}$};
\draw[dotted] (30,5) node[left] {\small $5$} -- (40,5) -- (40,0) node[below] {\small $40$};
\draw[dotted] (30,10) node[left] {\small $10$} -- (47.14,10) -- (47.14,0) node[below] {\small $Q_1$};
\draw[dotted] (30,12) node[left] {\small $12$} -- (50,12) -- (50,0) node[below=8pt] {\small $50$};
\draw[thick,red] (40,5) -- (50,12);
\fill (40,5) circle (2pt) node[left=2pt,above] {\scriptsize $A(40, 5)$};
\fill (47.14,10) circle (2pt) node[left] {\scriptsize $C(Q_1, 10)$};
\fill (50,12) circle (2pt) node[above] {\scriptsize $B(50, 12)$};
\end{tikzpicture}

\bigskip
\begin{tikzpicture}[x=0.22cm,y=0.15cm]
\draw[->] (50,0) -- (78,0) node[right] {\small ពិន្ទុ};
\draw[->] (50,0) -- (50,40) node[left] {$f^{\uparrow}$};
\draw[dotted] (50,24) node[left] {\small $24$} -- (60,24) -- (60,0) node[below] {\small $60$};
\draw[dotted] (50,30) node[left] {\small $30$} -- (65.45,30) -- (65.45,0) node[below] {\small $Q_3$};
\draw[dotted] (50,35) node[left] {\small $35$} -- (70,35) -- (70,0) node[below=8pt] {\small $70$};
\draw[thick,red] (60,24) -- (70,35);
\fill (60,24) circle (2pt) node[below right] {\scriptsize $A'(60, 24)$};
\fill (65.45,30) circle (2pt) node[left] {\scriptsize $C'(Q_3, 30)$};
\fill (70,35) circle (2pt) node[above] {\scriptsize $B'(70, 35)$};
\end{tikzpicture}
\end{center}
\end{example}

\begin{formulaBox}[ជាទូទៅ]
បើ $L$ ជាថ្នាក់គោលក្រោម, $f^{\uparrow}$ ជាប្រេកង់កើនមុនបន្ទាប់នៃថ្នាក់ដែល $Q$ ស្ថិតនៅ, $f$ ជាប្រេកង់, $i$ ជាប្រវែងថ្នាក់, $n$ ជាផលបូកប្រេកង់សរុប គេបាន ៖
\[
\begin{aligned}
Q_1 &= L + \frac{\frac{1}{4}n - f^{\uparrow}}{f}\times i\\
Q_3 &= L + \frac{\frac{3}{4}n - f^{\uparrow}}{f}\times i\\
L_P &= L + \frac{\frac{P}{100}n - f^{\uparrow}}{f}\times i
\end{aligned}
\]
\end{formulaBox}

\begin{example}[លំហាត់គំរូ ២]
ខាងក្រោមនេះជាតារាងបំណែងចែកប្រេកង់នៃទិន្នន័យ ការប្រាក់ (គិតជាដុល្លារ) របស់អ្នកវិនិយោគចំនួន $460$ នាក់ទទួលបានក្នុងមួយឆ្នាំ ។
\[
\begin{array}{|c|c|}\hline
\text{ការប្រាក់} & \text{ចំនួនអ្នកវិនិយោគ } f\\ \hline
30 - 40 & 17\\
40 - 50 & 55\\
50 - 60 & 142\\
60 - 70 & 153\\
70 - 80 & 93\\
\hline
\end{array}
\]
គណនា $Q_1$, $Q_3$ និង $L_{45}$ ។
\ans
ត្រូវបង្កើតតារាងប្រេកង់កើន ៖
\[
\begin{array}{|c|c|c|}\hline
\text{ការប្រាក់} & f & f^{\uparrow}\\ \hline
30 - 40 & 17 & 17\\
40 - 50 & 55 & 72\\
50 - 60 & 142 & 214\\
60 - 70 & 153 & 367\\
70 - 80 & 93 & 460\\ \hline
\end{array}
\]
\begin{itemize}
  \item $Q_1 = L + \dfrac{\frac{1}{4}n - f^{\uparrow}}{f}\times(i)$ ។ ដើម្បីអនុវត្តរូបមន្តនេះបាន លុះត្រាតែយើងដឹង $Q_1$ ស្ថិតក្នុងថ្នាក់មួយណាសិន ។ យើងដឹងថា $Q_1 =$ តម្លៃតួទី $\dfrac{1}{4}n =$ តម្លៃតួទី $\dfrac{1}{4}\times 460 =$ តម្លៃតួទី $115$ ។ តាមតារាងប្រេកង់កើន គេបាន $Q_1$ ស្ថិតក្នុងថ្នាក់ $50 - 60$ ។ គេបាន $L = 50$, $\dfrac{1}{4}n = 115$, $f^{\uparrow} = 72$, $f = 142$, $i = 60 - 50 = 10$ ។ ដូច្នេះ
  \[
  Q_1 = 50 + \frac{115 - 72}{142}\times 10 \approx 53.03 \text{ ដុល្លារ}
  \]
  \item $Q_3 =$ តម្លៃតួទី $\dfrac{3}{4}n =$ តម្លៃតួទី $\dfrac{3}{4}\times 460 =$ តម្លៃតួទី $345$ ។ តាមតារាងប្រេកង់កើន គេបាន $Q_3$ ស្ថិតក្នុងថ្នាក់ $60 - 70$ ។ គេបាន $L = 60$, $\dfrac{3}{4}n = 345$, $f^{\uparrow} = 214$, $f = 153$, $i = 70 - 60 = 10$ ។ ដូច្នេះ
  \[
  Q_3 = 60 + \frac{345 - 214}{153}\times 10 \approx 68.56 \text{ ដុល្លារ}
  \]
  \item $L_{45} =$ តម្លៃតួទី $\dfrac{45}{100}\times 460 =$ តម្លៃតួទី $207$ ។ តាមតារាងប្រេកង់កើន គេបាន $L_{45}$ ស្ថិតក្នុងថ្នាក់ $50 - 60$ ។ គេបាន $L = 50$, $\dfrac{P}{100}\times n = 207$, $f^{\uparrow} = 72$, $f = 142$, $i = 10$ ។ ដូច្នេះ
  \[
\begin{aligned}
L_{45} &= 50 + \frac{207 - 72}{142}\times 10 = 50 + 9.51\\
&= 59.51 \text{ ដុល្លារ}
\end{aligned}
\]
\end{itemize}
\textbf{បំណកស្រាយ ៖} ដូចជា $L_{45} = 59.51$ មានន័យថាមាន $45\%$ នៃអ្នកវិនិយោគទាំងអស់ទទួលបានការប្រាក់តិចជាង $59.51$ ដុល្លារក្នុងមួយឆ្នាំ ។
\end{example}

\begin{note}[សម្គាល់]
ចំពោះទិន្នន័យផ្ដុំជាថ្នាក់ការគណនាកាទីល យើងប្រើរូបមន្តរកទីតាំង $\dfrac{1}{4}n$, $\dfrac{3}{4}n$ ជំនួស $\dfrac{1}{4}(n + 1)$, $\dfrac{3}{4}(n + 1)$ ដែលជាទិន្នន័យមិនផ្ដុំជាថ្នាក់ ករណីនេះដូចគ្នាចំពោះភាគដប់ ឬភាគរយដែរ ។
\end{note}

\begin{pro}
តារាងបង្ហាញពីពិន្ទុសិស្សដែលបានធ្វើតេស្តភាសាអង់គ្លេសចំនួន $40$ នាក់ ។
\[
\begin{array}{|c|c|}\hline
\text{ពិន្ទុ} & \text{ចំនួនសិស្ស } f\\ \hline
50 - 60 & 8\\
60 - 70 & 10\\
70 - 80 & 13\\
80 - 90 & 6\\
90 - 100 & 3\\
\hline
\end{array}
\]
គណនា $L_{50}$, $Q_1$ និង $Q_3$ រួចបកស្រាយចម្លើយ ។
\ans
$f^{\uparrow}$ ៖ $8, 18, 31, 37, 40$ ។ $L_{50} = 70 + \dfrac{20 - 18}{13}\times 10 \approx 71.54$, $Q_1 = 60 + \dfrac{10 - 8}{10}\times 10 = 62$, $Q_3 = 70 + \dfrac{30 - 18}{13}\times 10 \approx 79.23$ ។ មានន័យថា $25\%$ នៃសិស្សបានពិន្ទុតិចជាង $62$, ពាក់កណ្ដាលបានតិចជាង $71.54$ និង $75\%$ បានតិចជាង $79.23$ ។
\end{pro}

\section*{លំហាត់អនុវត្តន៍}
\begin{enumerate}[1]
  \item គណនាកាទីលទាំងបីនិងភាគដប់ រួចបកស្រាយចម្លើយនៃសំណុំទិន្នន័យនីមួយៗខាងក្រោម ៖ \quad ក. $1, 3, 5, 5, 7, 0, 8, 2, 2, 4, 4$ \quad ខ. $7, 10, 22, 2, 0, 8, 12, 13, 18, 21$ ។
  \ans
  ក. រៀប ៖ $0, 1, 2, 2, 3, 4, 4, 5, 5, 7, 8$ ($n = 11$) ។ $Q_1 = 2$, $Q_2 = 4$, $Q_3 = 5$ ។ $D_1 = 0.5$, $D_2 = 1.5$, $D_3 = 2$, $D_4 = 2.5$, $D_5 = 4$, $D_6 = 4.5$, $D_7 = 5$, $D_8 = 6$, $D_9 = 7.5$ ។ \quad ខ. រៀប ៖ $0, 2, 7, 8, 10, 12, 13, 18, 21, 22$ ($n = 10$) ។ $Q_1 = 4.5$, $Q_2 = 11$, $Q_3 = 19.5$ ។ $D_1 = 1$, $D_2 = 4.5$, $D_3 = 7.5$, $D_4 = 9$, $D_5 = 11$, $D_6 = 12.5$, $D_7 = 15.5$, $D_8 = 19.5$, $D_9 = 21.5$ ។
  \item គណនាកាទីលទី $1$ $Q_1$ កាទីលទី $3$ $Q_3$ និង $L_{68}$ រួចបកស្រាយចម្លើយនៃសំណុំទិន្នន័យ ៖ \quad ក. $18, 0, 8, 16, 10, 4, 4, 14, 2, 13, 5, 7$ \quad ខ. $0.7, 0.4, 0.65, 0.78, 0.45, 0.32, 1.9, 0.0078$ ។
  \ans
  ក. រៀប ៖ $0, 2, 4, 4, 5, 7, 8, 10, 13, 14, 16, 18$ ។ $Q_1 =$ តួទី $3.25$ ៖ $4$, $Q_3 =$ តួទី $9.75$ ៖ $\dfrac{13 + 14}{2} = 13.5$, $L_{68} =$ តួទី $8.84$ ៖ $\dfrac{10 + 13}{2} = 11.5$ ។ \quad ខ. រៀប ៖ $0.0078, 0.32, 0.4, 0.45, 0.65, 0.7, 0.78, 1.9$ ។ $Q_1 =$ តួទី $2.25$ ៖ $0.36$, $Q_3 =$ តួទី $6.75$ ៖ $0.74$, $L_{68} =$ តួទី $6.12$ ៖ $0.74$ ។
  \item គណនាមេដ្យាន កាទីលទី $1$ $Q_1$ កាទីលទី $3$ $Q_3$ និង $L_{63}$ រួចបកស្រាយចម្លើយនៃបំណែងចែកប្រេកង់ខាងក្រោម ៖
  \[
  \text{ក.}\ \begin{array}{|c|cccccc|}\hline x & 4 & 5 & 6 & 7 & 8 & 9\\ \hline f & 12 & 9 & 8 & 13 & 9 & 9\\ \hline\end{array}
  \qquad
  \text{ខ.}\ \begin{array}{|c|ccccc|}\hline x & 12 & 13 & 14 & 15 & 16\\ \hline f & 3 & 9 & 11 & 15 & 7\\ \hline\end{array}
  \]
  \ans
  ក. $n = 60$, $f^{\uparrow}$ ៖ $12, 21, 29, 42, 51, 60$ ។ មេដ្យាន (តួទី $30.5$) $= 7$, $Q_1$ (តួទី $15.25$) $= 5$, $Q_3$ (តួទី $45.75$) $= 8$, $L_{63}$ (តួទី $38.43$) $= 7$ ។ \quad ខ. $n = 45$, $f^{\uparrow}$ ៖ $3, 12, 23, 38, 45$ ។ មេដ្យាន (តួទី $23$) $= 14$, $Q_1$ (តួទី $11.5$) $= 13$, $Q_3$ (តួទី $34.5$) $= 15$, $L_{63}$ (តួទី $28.98$) $= 15$ ។
  \item តារាងខាងក្រោមបង្ហាញពិន្ទុ ដែលបានមកពីបោះគ្រាប់ឡុកឡាក់មួយ ចំនួន $60$ ដង ។ គណនា $Q_2$ និង $D_1$ រួចបកស្រាយចម្លើយ ។
  \[
  \begin{array}{|l|cccccc|}\hline
  \text{ពិន្ទុ } x & 1 & 2 & 3 & 4 & 5 & 6\\ \hline
  \text{ប្រេកង់ } f & 12 & 9 & 8 & 13 & 9 & 9\\ \hline
  \end{array}
  \]
  \ans
  $f^{\uparrow}$ ៖ $12, 21, 29, 42, 51, 60$ ។ $Q_2 =$ តួទី $30.5$ ៖ $Q_2 = 4$ ។ $D_1 =$ តួទី $\dfrac{1}{10}\times 61 = 6.1$ ៖ $D_1 = 1$ ។ ពាក់កណ្ដាលនៃការបោះបានពិន្ទុមិនលើស $4$ ។
  \item អ្នកលក់សៀវភៅម្នាក់បានឱ្យដឹងថា សៀវភៅថ្មីរបស់គាត់មានទាំងអស់ចំនួន $25$ ក្បាល ក្នុងនោះមានចំនួន $4$ ក្បាលមួយក្បាលលក់តម្លៃ $\$10$ មានចំនួន $6$ ក្បាលមួយក្បាលលក់តម្លៃ $\$8$ មានចំនួន $4$ ក្បាលមួយក្បាលលក់តម្លៃ $\$7$ មានចំនួន $9$ ក្បាលមួយក្បាលលក់តម្លៃ $\$15$ និងមានចំនួន $2$ ក្បាលមួយក្បាលលក់តម្លៃ $\$22$ ។ \quad ក. គណនាតម្លៃមធ្យមក្នុងសៀវភៅមួយក្បាល \quad ខ. គណនាកាទីល $Q_1$, $Q_2$, $Q_3$ និង $L_{32}$ រួចបកស្រាយចម្លើយ ។
  \ans
  ក. $\bar{x} = \dfrac{4(10) + 6(8) + 4(7) + 9(15) + 2(22)}{25} = \dfrac{295}{25} = \$11.8$ ។ \quad ខ. រៀប ៖ $\$7(4)$, $\$8(6)$, $\$10(4)$, $\$15(9)$, $\$22(2)$; $f^{\uparrow}$ ៖ $4, 10, 14, 23, 25$ ។ $Q_1$ (តួទី $6.5$) $= \$8$, $Q_2$ (តួទី $13$) $= \$10$, $Q_3$ (តួទី $19.5$) $= \$15$, $L_{32}$ (តួទី $8.32$) $= \$8$ ។
  \item ការពិសោធន៍បណ្ដុះកូនប៉េងប៉ោះដោយប្រើជីគីមីក្នុងរយៈពេលមួយអាទិត្យ បម្រែបម្រួលកម្ពស់កូនប៉េងប៉ោះ ទទួលបានលទ្ធផលដូចខាងក្រោម ៖
  \[
  \begin{array}{|c|c|}\hline
  \text{កម្ពស់ (mm)} & \text{ចំនួនកូនប៉េងប៉ោះ } f\\ \hline
  0 - 5 & 3\\
  5 - 10 & 15\\
  10 - 15 & 72\\
  15 - 20 & 15\\
  20 - 25 & 90\\
  25 - 30 & 36\\
  30 - 35 & 8\\
  \hline
  \end{array}
  \]
  គណនា $Q_1$, $Q_3$ និង $D_1$ រួចបកស្រាយចម្លើយ ។
  \ans
  $n = 239$, $f^{\uparrow}$ ៖ $3, 18, 90, 105, 195, 231, 239$ ។ $Q_1$ ៖ $\dfrac{n}{4} = 59.75$ ក្នុងថ្នាក់ $10 - 15$ ៖ $Q_1 = 10 + \dfrac{59.75 - 18}{72}\times 5 \approx 12.90$ mm ។ $Q_3$ ៖ $\dfrac{3n}{4} = 179.25$ ក្នុងថ្នាក់ $20 - 25$ ៖ $Q_3 = 20 + \dfrac{179.25 - 105}{90}\times 5 \approx 24.13$ mm ។ $D_1$ ៖ $\dfrac{n}{10} = 23.9$ ៖ $D_1 = 10 + \dfrac{23.9 - 18}{72}\times 5 \approx 10.41$ mm ។
  \item គេបានកត់ត្រា ចំនួនសៀវភៅដែលមាននៅតាមធ្នើរ ក្នុងបណ្ណាល័យមួយ ។
  \[
  \begin{array}{|c|c|}\hline
  \text{ចំនួនសៀវភៅ} & \text{ចំនួនធ្នើរ}\\ \hline
  30 - 35 & 4\\
  35 - 40 & 6\\
  40 - 45 & 10\\
  45 - 50 & 13\\
  50 - 55 & 5\\
  55 - 60 & 3\\
  \hline
  \end{array}
  \]
  គណនាមេដ្យាន និង $L_{28}$ រួចបកស្រាយចម្លើយ ។
  \ans
  តាមតារាង ចំនួនធ្នើរសរុប $n = 41$, $f^{\uparrow}$ ៖ $4, 10, 20, 33, 38, 41$ ។ មេដ្យាន ៖ $\dfrac{n}{2} = 20.5$ ក្នុងថ្នាក់ $45 - 50$ ៖ $45 + \dfrac{20.5 - 20}{13}\times 5 \approx 45.19$ ។ $L_{28}$ ៖ $0.28\times 41 = 11.48$ ក្នុងថ្នាក់ $40 - 45$ ៖ $40 + \dfrac{11.48 - 10}{10}\times 5 = 40.74$ ។
\end{enumerate}
